\documentclass[french]{report} 
% valeur possibles : article, proc, book, report, letter, slides
% options possibles : twosides (recto-verso), 12pt, 11pt, ... (taille de police), draft (brouillon)

% pour l'edition de document en français:
\usepackage[utf8]{inputenc} % codage du fichier source
\usepackage[T1]{fontenc} % codage des fontes TeX
\usepackage{babel} % document en français

% Mise en page, propreté de rédaction
\usepackage[bottom=2cm,top=2cm,right=2cm,left=2cm]{geometry} % pour les marges
\usepackage{multicol} % pour ecrire sur plusieurs colones localement
\usepackage{xcolor} %metre du texte en couleur

% ecrire des maths
\usepackage{amsmath, amsfonts, amssymb}




%%%%%%% COMMANDES   
\newcommand{\N}[1]{\mathbb{N}^{#1}}

\newcommand{\alh}[1]{\textcolor{teal}{#1}}
\newcommand{\md}[1]{\textcolor{cyan}{#1}}



% Infos du document
\title{Interprétabilité pour l'Extraction d'Information dans les systèmes de RAG}
\author{Marine DELVALLEZ}
\date{\today}

\begin{document}
    \maketitle

    \chapter*{Remerciements}

    \tableofcontents

    \chapter{Introduction}
    \section{Environnement du stage}

    \begin{verbatim}
- Stage à lieu au LIFO, dans l'équipe CA, encadré par ...
- Stage de fin de master d'informatique Minerve GPEx
- Challenge@EvalLLM comme cas d'usage, participation envisagée mais timing pas OK
- différents événements et sorties organisées pendant le stage 
  - TransIA : Tours, Présentations pluridisciplinaires sur l'intégration de l'IA (sous toutes ses formes) 
  dans la société
  - TAL-I : journée du GDR TaL-I à Paris (Jussieu) sur les travaux portant sur l'interprétabilité et 
  l'explicabilité dans le contexte des tâches de traitement de la langue naturelle
  - 2nd Symposium on Mental Health and AI : deux jours de conférence (suivi en pointillé en distanciel) 
  présentation des travaux à la croisée entre santé mentale et intelligence artificielle (impact, outils 
  de détection, de prise en charge)
  - Place du Numérique : événement public de sensibilisation et de communication sur le Numérique 
  organisé place du Martoi à Orléans (Des événements similaires organisés dans chaque département 
  de la région CVL) J'ai participé à la tenue du stand de l'université, été ambassadrice de l'événement 
  et ai joué le rôle d'un témoin dan sle tribunal des génération futures portant sur "L'IA nous 
  contrôle-t-elle?
  - Présentation aux journées d'équipe CA
    \end{verbatim}


\section{Outils utilisés durant le stage}

    \begin{verbatim}
- Développement de l'architecture et réalisation des expérimentations en python dans 3 formats 
(script, implémentation d'un CLI, notebook)
- Principales librairies python utilisées : PyTorch, TransformerLens, HuggingFace, PyMuPDF4LLM,...
- Utilisation de plusieurs git pour assurer le suivi des versions des outils développé lors du stage
- Pour l'execution de tâches demandantes en charge de calcul (Génération de la base de représentation 
des données, execution des différentes architectures explorées et/ou implémentées) : utilisation des 
grappes de calcul accessibles aux chercheurs du LIFO CaSciModOT et Mirev
- utilisation d'outils à base d'IA très modérée et pour les tâches de l'ordre de reformulation
(ChatGPT), traduction (DeepL), exploration de la littérature (LitMaps, Perplexity), rédaction 
de petites fonctions et outils python (ChatGPT)
    \end{verbatim}



\section{Contexte et enjeux du sujet}

    \begin{verbatim}
TODO        
    \end{verbatim}


    \chapter{État de l'art}
    \section{Systèmes de Génération et Grands Modèles de Langue Augmentés par Extraction (RAG et RALLM)}

    \subsection{Besoin et Définition}
        \begin{verbatim}
- Les modèles profonds apprennent sur les données [ref à récupérer de Lewis+20]
- Le problème des modèles génératifs : hallucination et connaissance bornée dans 
le temps [ref à récupérer de Lewis+20]
- Il y a quelques années, on a cherché à associer un base de connaissances à ces 
modèles [Lewis+20] (= définition de RAG)
- Parmis ces approches, on retrouve le cas des modèles de langue : RALLM
        \end{verbatim}


    \subsection{L'architecture RAG}

        \begin{verbatim}
Description de l'archi de façon plus avancée
- Ces modèles peuvent être utilisés pour différents types de tâches (voir tableau 
csv sur RAGRights)
- Avec le temps, l'idée de [Lewis+20] a été adaptée et augmentée de composants 
supplémentaires
- [Fan+23] propose la structure générique ... (exploiter slides Séminaire CA)
        \end{verbatim}


    \subsection{Architectures aperçues au cours du stage}

        \begin{verbatim}
Au delà de l'innovation de l'architecture, on retrouve des innovations sur plusieurs 
caractéristiques de ces modèles : [R24/04] [R16/02]
- Multimodal [Liu+25]
- Chain of Thought [Xu+24]
- Gestion BDD dynamique [ref à trouver dans R24/05]
- BDD sparse et interprétable [Prouteau]???
        \end{verbatim}


    \subsection{Architecture et Modèles utilisées au cours du stage}

        \begin{verbatim}
Dans le cadre du stage, 
- travailler sur une adaptation du modèle RAG4HN. [Tran+25] 
- _description de l'architecture RAG4HN_ + parallèle avec l'architecture générique 
de [Fan+23]
- L'extracteur de RAG4HN est mE5 [Wang+25b] (version multilingue de E5 [Whang+25a])
        \end{verbatim}


\section{Systèmes RAG de confiance}

    \subsection{Trustworthy AI}

        \begin{verbatim}
La définition d'IA de confiance (Trustworthy AI) tend à varier en fonction des situation
et des utilisations de cette expression. Danss le contexte du Machine Learning, on désigne 
l'IA de confiance comme l'ensembles des systèmes d'IA dans lesquels il est légitime (dans 
le sens humainement entendable) d'avoir confiance. Le NIST (National Institute of Standards 
and Technology (US)) propose différents axes ou composantes de l'IA de confiance : ...
        \end{verbatim}


    \subsection{Explicabilité pour le RAG}
        \begin{verbatim}
Dans la suite, nous nous concentrons sur l'explicabilité dans le contexte du RAG. [Ni+25] 
distingue deux temps pour l'explicabilité : Extraction et génération ... (prévoir la question 
de l'association/intégration des deux)
- Métrique ??
        \end{verbatim}


    \subsection{Interprétabilité Mécaniste}
        \begin{verbatim}
[Somvanshi+25]
- Interprétabilité vs Explicabilité
- Définition d'Interprétabilité Mécaniste et approches (Manual Circuit Tracing,
*Intervention-based technique*, Representation analysis, Toy Model and synthetic tasks)
        \end{verbatim}


    \subsection{Activation Patching}

        \begin{verbatim}
- Définition de activation patching [Chen+24] [R30/04] [Séminaire CA]
- MechIR [Parry+25] 
- Perturbation
- Données produites, visualisation et interprétation des graphiques
        \end{verbatim}


    \chapter{Construction d'un dataset}
    \section{Données du Challenge RAG@EvalLLM}

    \begin{verbatim}
Présentation des données : 
- pdf à traiter sans structure normée du contenu
- majoritairement en français sauf quelques uns en anglais
- peuvent contenir des images
Présentation de la tâche du challenge (sauf si fait dans introduction)
- Identification des documents pertinents pour une question (phrase 
parfois complexe ou suite de mots clés) 
- questions fournies au dernier moment
    \end{verbatim}

\section{Structure de données visée}
    \begin{verbatim}
Base de données:        
- Paires 
- liste des champs et justification
Choix pour datasets
- Dataset avec des paires question-document (questions différentes et 
différents pdf, taille des documents, répartition positif/négatif)
- Dataset avec une seule question (construction à partir de 3 pdf) 
[description du dataset dans les expés et les specs]
    - motivation (reprendre l'expé présentée par MechIR, faciliter la création)
    - limites (biais, diversité ~> généralisabilité)
=> Placer le nom des outils utilisés (PyMuPDF4LLM et langchain.text_splitter.MarkdownSplitter)

    \end{verbatim}


    \chapter{Adaptation du modèle RAG4HN}
    \section{Architecture RAG4HistoricalNewspapers}

    \begin{verbatim}
Reprendre [Tran+25]
    \end{verbatim}

\section{Adaptation du modèle pour le stage}

    \begin{verbatim}
- Suppression de la partie WebSearch (hors sujet)
- Suppression de la première extraction sur la base des titres 
(incompatible avec le format des données du challenge)
- Passage de E5 à mE5 (on manipule du Français) [/!\ à vérifier]
- Montée en version des librairies pour la compatibilité avec les outils 
- Reprise du reranking (mettre au propre la photo du schema sur l'ardoise)
    \end{verbatim}


    \chapter{Cartographie d'un modèle d'extraction d'information}
    \section{Objectifs}

    \begin{verbatim}
- Objectif naif : Identifier les composants du modèle impliqués dans l'identification et la gestion 
de la langue de spécialité d'un dataset
- Objectif pratique : Identifier et proposer des préconisations d'usage de la librairie MechIR pour 
la localisation des composant impliqués dans la langue de spécialité
- Définition de langue de spécialité, réduction de l'objectif à l'identification du vocabulaire 
de spécialité 
    \end{verbatim}


\section{Méthodologie}
    \subsection{Construction des perturbations}
        \begin{verbatim}
- on s'intéresse aux perturbations de type replace pour travailler sur l'impact des mots importants
(car liés au vocabulaire de spécialité) en les replaçant
- On propose plusieurs candidats au remplacement classé en 4 niveaux ...
- Le choix des remplaçant est fait par jugement humain
- Le choix des mots à remplacer se fait à l'aide de 
    - fréquence dans l'échantillon de données sélectionné
    - IDF 
    - appartenance au vocabulaire de spécialité (jugement humain)
    - présence du mot dans la question étudiée (cas du dataset à 1! question)
- pour augmenter le nombre de document perturbé et limiter l'impact d'un mot seulement dans la 
cartographie, on propose un nouveau type de perturbation : full_replace
        \end{verbatim}

    \subsection{Estimation de la pertinence d'une perturbation}
        \begin{verbatim}
SI PAS DÉJÀ PRÉSENTÉ DANS PARTIE DE L'ÉTAT DE L'ART SUR MECHIR
- On visualise des scores de similarité entre questoin et document
- graphique d'exemple + texte d'accompagnement à la lecture
- Un perturbation peut être considérée comme pertinente lorsqu'elle modifie significativement 
(jugement humain) la distribution des cores pour l'échantillon du dataset étudié
        \end{verbatim}

    \subsection{Interprétation des résultats de l'Activation Patching}

        \begin{verbatim}
SI PAS DÉJÀ FAITE DANS SOTA sur MechIR
- graphique + accompagnement à la lecture (on obtient une sensibilité de chaque 
noeuds pour chaque paire (Question, document), on visualise la moyenne)

- étant donné la construction de ces valeurs, on s'intéresse aussi aux écarts-types
- pour faciliter la visualisation et la comparaison en cas de sensibilité variable 
d'une perturbation à l'autre on propose de rescaler les poids des matrices
- L'identification des noeuds sensibles doit reposer sur un critère commun. On propose
un critère de sensibilité ...
- Visualisation des noeuds sensibles commun à plusieurs perturbation (matrice des noeuds sensibles)
    -> croiser les différentes caractéristiques (mot remplacé, type, niveau)
    -> critique des noeuds sensibles de la dernière couche
        \end{verbatim}

\section{Résultats et Analyse}

    \subsection{Application à la cartographie du vocabulaire de la défense français dans mE5}
        \begin{verbatim}
- Grand nombre de noeuds de la dernière couche du modèle
- pas de noeuds clairement associé au vocabulaire
...
        \end{verbatim}


    \subsection{Préconisations d'utilisation}
        \begin{verbatim}
- des perturbations qui couvrent tous ou très grande partie des documents étudiés
- utiliser la question comme document témoin
        \end{verbatim}


\section{Perspectives d'évolution de la méthode proposée}

    \begin{verbatim}
- échelle pour situer le niveau d'un mot candidat remplaçant en fonction du mot remplacé
- critère d'identification du vocabulaire de spécialité
- intégration d'autres informations au critère de sensibilité
    \end{verbatim}

    \chapter{Conclusions et Perspectives}
    \section{Conclusion}

    Au cours de ce stage, j'ai exploré la littérature scientifique portant sur le RAG et son explicabilité. J'ai découvert les différentes perspectives d'évolution de ces systèmes et les perspectives d'explicabilité pour les tâches l'extraction d'information à l'aide des méthodes d'interprétabilité mécaniste. J'ai repris le système RAG proposé par \cite{tran_retrieval_2024} et conçu un dataset type à partir des données fournies par le challenge RAG@EvalLLM. Enfin, j'ai proposé une méthodologie d'usage de l'activation patching avec l'outil MechIR \cite{parry_mechir_2025}. Cette dernière inclus la conception des perturbations ainsi que la visualisation et l'interprétation des résultats fournis par MechIR.

\section{Perspectives générales}

    Au cours du stage, j'ai finalement choisi de poursuivre un sujet de thèse différent de celui visé par ce stage. Les travaux réalisés  constituent néanmoins une première exploration méthodologique autour de l'interprétabilité mécaniste appliquée aux modèles d'extraction d'information pour le RAG.

    Dans l'éventualité d'une reprise de ces travaux, voici quelques perspectives pour tende vers le résultat initialement prévu. Dans un premier temps, L'identification des nœuds proposée ne constitue pas encore une cartographie d'un modèle d'extraction. C'est la première étape. Une fois la cartographie obtenue, l'extraction d'explication du fonctionnement du modèle lors de l'usage doit être définie. Il est nécessaire de garder en tête que l'utilisateur visé n'est pas informaticien. Par conséquent, un travail doit être effectué pour rendre les explications accessible au public visé.

    Par ailleurs, d'autres perspectives d'usage d'une telle cartographie existent. Cette dernière peut servir d'orientation pour améliorer ou spécialiser un modèle. De premiers résultats existent déjà sur l'utilisation de l'activation patching appliqué au prunning.

\end{document}